% SEGMENT OLP-0060-S01
% Part: first-order-logic
% Chapter: syntax-and-semantics
% Section: formation-sequences

\documentclass[../../../include/open-logic-section]{subfiles}

\begin{document}

\olfileid{pl}{syn}{fseq}

\olsection{அமைப்புத் தொடர்முறைகள்}

!!{formula}s ஐத் தொகுத்தறிதல் வரையறையால் வரையறுப்பதும், அதற்கு
நிறைவாக !!{formula}s இன் பண்புகளைத் தொகுத்தறிதலால் நிறுவுவதும் நேர்த்தியான,
செயல்திறமிக்க அணுகுமுறையாகும். எனினும், !!{formula}s இன் கட்டமைப்பை
அடிமட்டத்திலிருந்து மேல்நோக்கிப் படிப்படியாக அணுகுவதும் பயனுள்ளதாக
இருக்கலாம். இங்கு இதனை \emph{அமைப்புத் தொடர்முறை} என்ற கருத்தைப்
பயன்படுத்திச் செய்கிறோம்.

\begin{defn}[வாய்பாடுகளுக்கான அமைப்புத் தொடர்முறைகள்]
\ollabel{defn:fseq-frm}
$\Lang L_0$ மொழியின் குறியீட்டுச் சரங்களாலான முடிவுறு தொடர்முறை
$\tuple{!A_0,\dotsc,!A_n}$, $!A$ க்கான ஓர் \emph{அமைப்புத்
தொடர்முறை} என்பதற்கான நிபந்தனைகள்: $!A \ident !A_n$; மேலும் ஒவ்வொரு
$i \leq n$ க்கும், $!A_i$ ஓர் அணு வாய்பாடாக இருக்க வேண்டும் அல்லது
பின்வருவனவற்றுள் ஒன்று பொருந்துமாறு $j,k < i$ இருக்க வேண்டும்:
\begin{enumerate}
    \tagitem{prvNot}{$!A_i \ident \lnot !A_j$.}{}%
    \tagitem{prvAnd}{$!A_i \ident (!A_j \land !A_k)$.}{}%
    \tagitem{prvOr}{$!A_i \ident (!A_j \lor !A_k)$.}{}%
    \tagitem{prvIf}{$!A_i \ident (!A_j \lif !A_k)$.}{}%
    \tagitem{prvIff}{$!A_i \ident (!A_j \liff !A_k)$.}{}%
\end{enumerate}
\end{defn}

% SEGMENT OLP-0060-S02
\begin{ex}
\[
\tuple{
    \Obj p_0,
    \Obj p_1,
    (\Obj p_1 \land \Obj p_0),
    \lnot (\Obj p_1 \land \Obj p_0)
}
\]
என்பது $\lnot (\Obj p_1 \land \Obj p_0)$ இன் ஓர் அமைப்புத்
தொடர்முறை; பின்வருவதும் அத்தகையதே:
\[
\tuple{
    \Obj p_0,
    \Obj p_1,
    \Obj p_0,
    (\Obj p_1 \land \Obj p_0),
    (\Obj p_0 \lif \Obj p_1),
    \lnot (\Obj p_1 \land \Obj p_0)
}.
\]
%
இரண்டாம் எடுத்துக்காட்டிலிருந்து காணக்கூடியதுபோல், அமைப்புத்
தொடர்முறைகளில் ``தேவையற்றவை'' இருக்கலாம்: மிகையான அல்லது கட்டமைப்புக்குப்
பங்களிக்காத வாய்பாடுகள் இவ்வாறு அமையலாம்.
\end{ex}

% SEGMENT OLP-0060-S03
\begin{prop}
\ollabel{prop:formed}
$\Frm[L_0]$ இல் உள்ள ஒவ்வொரு !!{formula}~$!A$ உம் ஓர் அமைப்புத்
தொடர்முறையைக் கொண்டுள்ளது.
\end{prop}

\begin{proof}
$!A$ அணு வாய்பாடு எனக் கொள்க. அப்போது தொடர்முறை~$\tuple{!A}$ ஆனது
$!A$ க்கான அமைப்புத் தொடர்முறை.
%
இப்போது $!B$, $!C$ ஆகியவை முறையே $\tuple{!B_0,\dotsc,!B_n}$,
$\tuple{!C_0,\dotsc,!C_m}$ என்ற அமைப்புத் தொடர்முறைகளைக் கொண்டுள்ளதாகக்
கொள்க.
%
\begin{enumerate}
  \tagitem{prvNot}{$!A \ident \lnot !B$ எனில்,
      $\tuple{!B_0,\dotsc,!B_n,\lnot !B_n}$ ஆனது $!A$ க்கான அமைப்புத்
      தொடர்முறை.}{}
  \tagitem{prvAnd}{$!A \ident (!B \land !C)$ எனில்,
      $\tuple{!B_0,\dotsc,!B_n,!C_0,\dotsc,!C_m,(!B_n \land !C_m)}$
      ஆனது $!A$ க்கான அமைப்புத் தொடர்முறை.}{}
  \tagitem{prvOr}{$!A \ident (!B \lor !C)$ எனில்,
      $\tuple{!B_0,\dotsc,!B_n,!C_0,\dotsc,!C_m,(!B_n \lor !C_m)}$
      ஆனது $!A$ க்கான அமைப்புத் தொடர்முறை.}{}
  \tagitem{prvIf}{$!A \ident (!B \lif !C)$ எனில்,
      $\tuple{!B_0,\dotsc,!B_n,!C_0,\dotsc,!C_m,(!B_n \lif !C_m)}$
      ஆனது $!A$ க்கான அமைப்புத் தொடர்முறை.}{}
  \tagitem{prvIff}{$!A \ident (!B \liff !C)$ எனில்,
      $\tuple{!B_0,\dotsc,!B_n,!C_0,\dotsc,!C_m,(!B_n \liff !C_m)}$
      ஆனது $!A$ க்கான அமைப்புத் தொடர்முறை.}{}
\end{enumerate}
!!{formula}s இன்மீதான தொகுத்தறிதல் நெறியின்படி, ஒவ்வொரு !!{formula}வும்
ஓர் அமைப்புத் தொடர்முறையைக் கொண்டுள்ளது.
\end{proof}

% SEGMENT OLP-0060-S04
மறுதிசையையும் நிறுவலாம். வாய்பாடுகளை வரையறுக்கும் நமது இரு முறைகளும்
சமானமானவை, அதாவது அவை ஒரே விளைவுகளைத் தருகின்றன, என்பதை இது காட்டுவதால்
இது முக்கியமானது. அமைப்புத் தொடர்முறைகளின் நீளத்தின்மீதான வழக்கமான
தொகுத்தறிதலைப் பயன்படுத்தி வாய்பாடுகளைப் பற்றிய தேற்றங்களை நிறுவலாம்
என்பதும் இதன் பொருள்.

\begin{lem}
\ollabel{lem:fseq-init}
$\tuple{!A_0,\dotsc,!A_n}$ ஆனது $!A_n$ க்கான ஓர் அமைப்புத்
தொடர்முறையாகவும், $k \leq n$ ஆகவும் இருக்கட்டும். அப்போது
$\tuple{!A_0,\dotsc,!A_k}$ ஆனது $!A_k$ க்கான ஓர் அமைப்புத் தொடர்முறை.
\end{lem}

\begin{proof}
பயிற்சி.
\end{proof}

% SEGMENT OLP-0060-S05
\begin{thm}
\ollabel{thm:fseq-frm-equiv}
$\Frm[L_0]$ என்பது அமைப்புத் தொடர்முறையைக் கொண்ட $\Lang L_0$ மொழியின்
எல்லா குறியீட்டுச் சரங்களின் கணமாகும்.
\end{thm}

% SEGMENT OLP-0060-S06
\begin{proof}
$\Lang L_0$ மொழியில் அமைப்புத் தொடர்முறையைக் கொண்ட எல்லா குறியீட்டுச்
சரங்களின் கணம் $F$ என்க. \olref[pl][syn][fseq]{prop:formed} இல்
$\Frm[L_0] \subseteq F$ என்று கண்டுள்ளோம்; இப்போது மறுதிசையை நிறுவுகிறோம்.

$!A$ ஆனது $\tuple{!A_0,\dotsc,!A_n}$ என்ற அமைப்புத் தொடர்முறையைக்
கொண்டிருக்கட்டும். $n$ இன்மீதான வலுத் தொகுத்தறிதலால்
$!A \in \Frm[L_0]$ என்று நிறுவுகிறோம். $m < n$ நீளமுள்ள அமைப்புத்
தொடர்முறையைக் கொண்ட ஒவ்வொரு குறியீட்டுச் சரமும் $\Frm[L_0]$ இல் உள்ளது
என்பதே நமது தொகுத்தறிதல் கருதுகோள். அமைப்புத் தொடர்முறையின்
வரையறையின்படி, $!A_n$ அணு வாய்பாடாக இருக்க வேண்டும் அல்லது
பின்வருவனவற்றுள் ஒன்று பொருந்துமாறு $j,k < n$ இருக்க வேண்டும்:
\begin{enumerate}
    \tagitem{prvNot}{$!A_n \ident \lnot !A_j$.}{}%
    \tagitem{prvAnd}{$!A_n \ident (!A_j \land !A_k)$.}{}%
    \tagitem{prvOr}{$!A_n \ident (!A_j \lor !A_k)$.}{}%
    \tagitem{prvIf}{$!A_n \ident (!A_j \lif !A_k)$.}{}%
    \tagitem{prvIff}{$!A_n \ident (!A_j \liff !A_k)$.}{}%
\end{enumerate}

% SEGMENT OLP-0060-S07
இப்போது நேர்வுகளால் காரணங்காட்டுவோம். $!A_n$ அணு வாய்பாடு எனில்,
$!A_n \in \Frm[L_0]$. மாறாக $!A \equiv (!A_j \land !A_k)$ எனக்
கொள்க. \olref[pl][syn][fseq]{lem:fseq-init} இன்படி,
$\tuple{!A_0,\dotsc,!A_j}$, $\tuple{!A_0,\dotsc,!A_k}$ ஆகியவை முறையே
$!A_j$, $!A_k$ க்கான அமைப்புத் தொடர்முறைகள். இவை $!A$ க்கான அமைப்புத்
தொடர்முறையின் மெய்யான தொடக்கத் துணைத்தொடர்முறைகள் என்பதால், இரண்டின்
நீளமும் $n$ ஐவிடக் குறைவு. எனவே தொகுத்தறிதல் கருதுகோளின்படி $!A_j$,
$!A_k$ ஆகியவை $\Frm[L_0]$ இல் உள்ளன. ஆகவே !!a{formula}வின்
வரையறையின்படி $(!A_j \land !A_k)$ உம் அதில் உள்ளது. மற்ற நேர்வுகளும்
இணையான காரணங்காட்டலைப் பின்பற்றுகின்றன.
\end{proof}

\end{document}
